% ======================================================
% Compile: xelatex Fe_H_derivation.tex
% Style: mimic IEEE Trans. Industrial Electronics (Meng 2025)
% Path:  reference/eq17_analysis/derivation/Fe_H_derivation.tex
% Note:  The complete F_e Row-3 form here (cols (3,5)/(3,7) nonzero via
%        F_1/F_2 control history) SUPERSEDES the truncated form in
%        derivation/phase1_Fe_derivation.md S10.1. The matching controller
%        code (6-arg build_F_e) is pending port from the Mac branch
%        learn/variance-walkthrough (commit 2aabec4); until that lands,
%        eq17_core.m still implements the truncated Row 3 (cols 5/7 = 0).
% ======================================================
\documentclass[11pt,a4paper]{ctexart}

\usepackage[fleqn]{amsmath}
\usepackage{amssymb,bm}
\usepackage{empheq}
\setlength{\mathindent}{0pt}
\usepackage{geometry}
\geometry{margin=2.2cm}
\usepackage{enumitem}
\usepackage{array,booktabs}
\usepackage{titlesec}
\usepackage{hyperref}
\hypersetup{colorlinks=true,linkcolor=blue,citecolor=blue}

% --- Tight layout (mimic IEEE Trans density) ---
\setlength{\parindent}{0pt}
\setlength{\parskip}{0.15em}
\setlength{\abovedisplayskip}{4pt}
\setlength{\belowdisplayskip}{4pt}
\setlength{\abovedisplayshortskip}{2pt}
\setlength{\belowdisplayshortskip}{2pt}
\titlespacing*{\section}{0pt}{1.4ex}{0.6ex}
\titlespacing*{\subsection}{0pt}{0.8ex}{0.3ex}

% ---------- Shortcuts ----------
\newcommand{\dt}{\Delta t}
\newcommand{\lc}{\lambda_c}
\newcommand{\fdx}{f_{dx}}
\newcommand{\fT}{f_T}
\newcommand{\dx}{\delta x}
\newcommand{\dxm}{\delta x_m}
\newcommand{\axm}{a_{xm}}
\newcommand{\dxone}{\delta x_1}
\newcommand{\dxtwo}{\delta x_2}
\newcommand{\dxthree}{\delta x_3}
\newcommand{\xD}{x_D}
\newcommand{\dxD}{\delta x_D}
\newcommand{\ax}{a_x}
\newcommand{\dax}{\delta a_x}
\newcommand{\axhat}{\hat a_x}
\newcommand{\xDhat}{\hat x_D}
\newcommand{\dxhat}{\delta\hat x}

% Bold for matrices / vectors
\newcommand{\Fe}{\mathbf{F}_e}
\newcommand{\HH}{\mathbf{H}}
\newcommand{\LL}{\mathbf{L}}
\newcommand{\xx}{\mathbf{x}}
\newcommand{\xxhat}{\hat{\mathbf{x}}}
\newcommand{\ee}{\mathbf{e}}
\newcommand{\qq}{\mathbf{q}}

% Estimation errors and composites
\newcommand{\ea}{e_{a}}
\newcommand{\exD}{e_{x_D}}
\newcommand{\dea}{\delta e_{a}}
\newcommand{\dexD}{\delta e_{x_D}}
\newcommand{\epsT}{\varepsilon_T}
\newcommand{\epsa}{\varepsilon_a}
\newcommand{\epsxD}{\varepsilon_{x_D}}
\newcommand{\etaR}{\eta_{R}}
\newcommand{\QQ}{\mathbf{Q}}

% Hat-state shortcuts
\newcommand{\dxonehat}{\delta\hat x_1}
\newcommand{\dxtwohat}{\delta\hat x_2}
\newcommand{\dxthreehat}{\delta\hat x_3}
\newcommand{\dxDhat}{\delta\hat x_D}
\newcommand{\daxhat}{\delta\hat a_x}

\begin{document}
\thispagestyle{empty}

\begin{center}
{\Large\bfseries Derivation of $\fdx[k]$, $\Fe[k]$, and $\HH$}
\end{center}
\vspace{0.3cm}

%=====================================================================
\section{Setup}
%=====================================================================

\subsection{Plant Equation}

\begin{align}
\gamma\,\dot x(t) &= \fdx(t) + \fT(t), \label{eq:plant-cont} \\
x[k{+}1] &= x[k] + \ax[k]\bigl\{\fdx[k]+\fT[k]\bigr\} + \xD[k]. \label{eq:plant-disc}
\end{align}

\subsection{Measurement}

For a sensor delay of $d=2$ steps,
\begin{equation}
\dxm[k] = \dx[k{-}d], \qquad \dx[k] := x_d[k] - x[k].
\label{eq:meas-dxm}
\end{equation}

%=====================================================================
\section{Derivation of the Motion Control Law}
\label{sec:ctrl-deriv}
%=====================================================================

This section constructs $\fdx[k]$ from \eqref{eq:plant-disc} so that the closed-loop tracking error follows a first-order pole $\lc \in [0,1)$. Forced approximations (no alternative) are labelled F1--F6 inline; their residual destinations are summarized in Appendix~\ref{app:catalog}.

\subsection{Design Objective}

The requirement is
\begin{equation}
\dx[k{+}1] \;\overset{!}{=}\; \lc\,\dx[k].
\label{eq:dctl-target}
\end{equation}
Substituting \eqref{eq:plant-disc} into \eqref{eq:dctl-target} and equating to $\lc\,\dx[k] = \lc x_d[k] - \lc x[k]$,
\begin{equation}
\ax[k]\bigl\{\fdx[k]+\fT[k]\bigr\}
   = x_d[k{+}1] - \lc\,x_d[k] - (1{-}\lc)\,x[k] - \xD[k].
\label{eq:dctl-balance}
\end{equation}

\subsection{Idealized Control Law}

If all quantities on the right of \eqref{eq:dctl-balance} are instantaneously available, solving for $\fdx$ yields
\begin{equation}
\fdx[k] = \ax^{-1}[k]\bigl\{x_d[k{+}1] - \lc\,x_d[k] - (1{-}\lc)\,x[k]\bigr\}
         - \fT[k] - \ax^{-1}[k]\,\xD[k].
\label{eq:dctl-ideal}
\end{equation}
The unknowns are $\ax[k]$, $x[k]$, $\fT[k]$, $\xD[k]$. The remaining derivation makes \eqref{eq:dctl-ideal} implementable using $\dxm[k]$, the past control history $\fdx[k{-}i]$, and the EKF estimates $\axhat[k{-}j]$, $\xDhat[k{-}j]$ for $j=0,\ldots,d$.

\subsection{State Reconstruction}

The sensor provides $x_m[k] = x[k{-}d] + n_x[k]$, not $x[k]$. Iterating \eqref{eq:plant-disc} from $k{-}d$ to $k$ (for $d=2$, two successive substitutions),
\begin{align}
x[k{-}1] &= x[k{-}2] + \ax[k{-}2]\{\fdx[k{-}2]+\fT[k{-}2]\} + \xD[k{-}2], \\
x[k]     &= x[k{-}1] + \ax[k{-}1]\{\fdx[k{-}1]+\fT[k{-}1]\} + \xD[k{-}1],
\end{align}
and summing the two yields the exact identity
\begin{equation}
x[k] = x[k{-}d] + \sum_{i=1}^{d}\!\bigl\{\ax[k{-}i]\bigl(\fdx[k{-}i]+\fT[k{-}i]\bigr) + \xD[k{-}i]\bigr\}.
\label{eq:tele-x}
\end{equation}

\subsection{Implementation Substitutions}

Each unknown in \eqref{eq:dctl-ideal} and \eqref{eq:tele-x} is replaced by an available quantity:
\begin{itemize}\setlength{\itemsep}{1pt}
  \item \textbf{F1} \quad $\ax[k] \mapsto \axhat[k]$ (EKF posterior).
  \item \textbf{F2} \quad $\ax[k{-}i] \mapsto \axhat[k{-}i]$, $i=1,\ldots,d$ (buffered past posteriors).
  \item \textbf{F3} \quad $x[k{-}d] \mapsto x_m[k]$ (sensor; the discarded $-n_x[k]$ reappears in the closed-loop noise composite).
  \item \textbf{F4} \quad $\fT[k]$ and $\fT[k{-}i] \mapsto 0$ (zero-mean white thermal force; unpredictable).
  \item \textbf{F5} \quad $\xD[k] \mapsto \xDhat[k]$ (EKF posterior).
  \item \textbf{F6} \quad $\xD[k{-}i] \mapsto \xDhat[k{-}i]$, $i=1,\ldots,d$ (buffered past posteriors).
\end{itemize}
After F2, F4, F6, the implementable reconstruction \eqref{eq:tele-x} becomes
\begin{equation}
\tilde x[k] = x_m[k] + \sum_{i=1}^{d}\!\bigl\{\axhat[k{-}i]\,\fdx[k{-}i] + \xDhat[k{-}i]\bigr\}.
\label{eq:xk-impl}
\end{equation}

\subsection{Final Control Law}

Substituting $\axhat[k]$, $\xDhat[k]$, and $\tilde x[k]$ into \eqref{eq:dctl-ideal}, dropping $\fT[k]$ per F4, and converting $x_m[k] = x_d[k{-}d] - \dxm[k]$,
\begin{empheq}[box=\fbox]{align}
\fdx[k] = \,& \axhat^{-1}[k]\bigl\{x_d[k{+}1] - \lc\,x_d[k] - (1{-}\lc)\,x_d[k{-}d] + (1{-}\lc)\,\dxm[k]\bigr\} \notag \\
           & - \axhat^{-1}[k]\,(1{-}\lc)\!\sum_{i=1}^{d}\!\bigl\{\axhat[k{-}i]\,\fdx[k{-}i] + \xDhat[k{-}i]\bigr\} \notag \\
           & - \axhat^{-1}[k]\,\xDhat[k].
\label{eq:ctrl-law}
\end{empheq}
The controller buffers $\{\axhat[k{-}i]\}_{i=1}^{d}$ and $\{\xDhat[k{-}i]\}_{i=1}^{d}$ in addition to its own past output $\{\fdx[k{-}i]\}_{i=1}^{d}$. Two design alternatives are discussed in the Appendix: the legacy passive form (\S\ref{app:passive}) that collapses the buffered past estimates, and the Meng~2023 limit (\S\ref{app:meng}) in which $\ax$ is calibrated and $\xD\equiv 0$.

%=====================================================================
\section{Closed-Loop Tracking Error Dynamics}
\label{sec:trk}
%=====================================================================

Define the estimation errors
\begin{equation}
\ea[k] := \ax[k] - \axhat[k], \qquad
\exD[k] := \xD[k] - \xDhat[k].
\label{eq:est-err-def}
\end{equation}

Since $\ax[k] = \axhat[k] + \ea[k]$, multiplying \eqref{eq:ctrl-law} by $\axhat[k]$ gives
\begin{align}
\ax[k]\,\fdx[k] = \,& \ea[k]\,\fdx[k] + x_d[k{+}1] - \lc\,x_d[k] - (1{-}\lc)\,x_d[k{-}d] + (1{-}\lc)\,\dxm[k] \notag\\
& - (1{-}\lc)\!\sum_{i=1}^{d}\!\bigl\{\axhat[k{-}i]\,\fdx[k{-}i] + \xDhat[k{-}i]\bigr\} - \xDhat[k].
\label{eq:axfdx}
\end{align}

Substituting \eqref{eq:plant-disc} and \eqref{eq:axfdx} into $\dx[k{+}1] = x_d[k{+}1] - x[k{+}1]$,
\begin{align}
\dx[k{+}1] = \,& \lc\,x_d[k] + (1{-}\lc)\,x_d[k{-}d] - x[k] - (1{-}\lc)\,\dxm[k] \notag\\
& + (1{-}\lc)\!\sum_{i=1}^{d}\!\bigl\{\axhat[k{-}i]\,\fdx[k{-}i] + \xDhat[k{-}i]\bigr\} \notag\\
& - \ea[k]\,\fdx[k] - \exD[k] - \ax[k]\,\fT[k].
\label{eq:substitute}
\end{align}

Substituting \eqref{eq:meas-dxm} for $\dxm[k]$ and \eqref{eq:tele-x} for $x[k]$, the $x_d[k{-}d]$ terms cancel; the remaining $\lc\bigl(x_d[k]-x[k{-}d]\bigr)$ returns to $\lc\,\dx[k]$ through \eqref{eq:tele-x}:
\begin{align}
\dx[k{+}1] = \,& \lc\,\dx[k] - \ea[k]\,\fdx[k] - \exD[k] - \ax[k]\,\fT[k] \notag\\
& - (1{-}\lc)\!\sum_{i=1}^{d}\!\bigl(\ax[k{-}i] - \axhat[k{-}i]\bigr)\fdx[k{-}i] \notag\\
& - (1{-}\lc)\!\sum_{i=1}^{d}\!\bigl(\xD[k{-}i] - \xDhat[k{-}i]\bigr)
  - (1{-}\lc)\!\sum_{i=1}^{d}\!\ax[k{-}i]\,\fT[k{-}i].
\label{eq:collect}
\end{align}

Applying \eqref{eq:est-err-def} at steps $k{-}i$,
\begin{align}
\dx[k{+}1] = \,& \lc\,\dx[k] - \ea[k]\,\fdx[k] - \exD[k] \notag\\
& - (1{-}\lc)\!\sum_{i=1}^{d}\!\bigl\{\ea[k{-}i]\,\fdx[k{-}i] + \exD[k{-}i]\bigr\} - \epsT[k],
\label{eq:closed-loop}
\end{align}
where
\begin{equation}
\epsT[k] := \ax[k]\,\fT[k] + (1{-}\lc)\!\sum_{i=1}^{d}\!\ax[k{-}i]\,\fT[k{-}i].
\label{eq:eps-T}
\end{equation}

Every residual in \eqref{eq:closed-loop} is an estimation error ($\ea$, $\exD$) or a thermal forcing ($\fT$); no Taylor expansion of $\ax$ is involved, hence no model-truncation residual appears.

For brevity, define
\begin{equation}
\mathcal{F}_1[k] := \sum_{i=1}^{d}\!\fdx[k{-}i], \qquad
\mathcal{F}_2[k] := \sum_{i=1}^{d}\!i\,\fdx[k{-}i].
\label{eq:F12}
\end{equation}

%=====================================================================
\section{State Vector and Output Matrix}
%=====================================================================

Estimation of $\axhat[k]$, $\dxhat[k]$, and $\xDhat[k]$ is required to update \eqref{eq:ctrl-law-passive}. Due to the $d$-step measurement delay, three variables $\dxone[k] = \dx[k{-}2]$, $\dxtwo[k] = \dx[k{-}1]$, $\dxthree[k] = \dx[k]$ model the motion error, and $\xD$ and $\ax$ are modeled as second-order processes. The augmented state vector is
\begin{equation}
\xx[k] = \bigl[\,\dxone,\ \dxtwo,\ \dxthree,\ \xD,\ \dxD,\ \ax,\ \dax\,\bigr]^T[k],
\label{eq:state-vec}
\end{equation}
whose evolution is given by
\begin{equation}
\begin{cases}
\dxone[k{+}1]   = \dxtwo[k] \\[2pt]
\dxtwo[k{+}1]   = \dxthree[k] \\[2pt]
\dxthree[k{+}1] = \lc\,\dxthree[k] - \fdx[k]\,\ea[k] - \exD[k] - (1{-}\lc)\!\sum_{i=1}^{d}\xD[k{-}i] - \epsT[k] + \etaR[k] \\[2pt]
\xD[k{+}1]      = \xD[k] + \dxD[k] \\[2pt]
\dxD[k{+}1]     = \dxD[k] \\[2pt]
\ax[k{+}1]      = \ax[k] + \dax[k] \\[2pt]
\dax[k{+}1]     = \dax[k]
\end{cases}
\label{eq:state-evo}
\end{equation}

The two measurements $y_1[k] = \dxm[k]$ and $y_2[k] = \axm[k]$ both inherit the $d$-step delay; in particular $y_1[k] = \dxone[k]$ and $\axm[k] \approx \ax[k{-}d] + n_a[k]$, where $n_a[k]$ is the intrinsic measurement noise. By the second-order process for $\ax$ in \eqref{eq:state-evo},
\begin{equation}
\ax[k{-}d] = \ax[k] - d\,\dax[k],
\label{eq:ax-back}
\end{equation}
so that $\axm[k] \approx \ax[k] - d\,\dax[k] + n_a[k]$, yielding
\begin{equation}
\HH = \begin{bmatrix}
1 & 0 & 0 & 0 & 0 & 0 & 0 \\
0 & 0 & 0 & 0 & 0 & 1 & -d
\end{bmatrix}.
\label{eq:H}
\end{equation}

%=====================================================================
\section{Closed-Loop Estimator and Error Dynamics}
\label{sec:est}
%=====================================================================

Through the feedback of two estimation errors $e_x[k] = \dxm[k] - \dxonehat[k]$ and $e_{a_x}[k] = \axm[k] - \axhat[k]$, the closed-loop estimator is designed as
\begin{equation}
\left\{\begin{aligned}
\dxonehat[k{+}1]   &= \dxtwohat[k] + \ell_{11}[k]\,e_x[k] + \ell_{12}[k]\,e_{a_x}[k] \\
\dxtwohat[k{+}1]   &= \dxthreehat[k] + \ell_{21}[k]\,e_x[k] + \ell_{22}[k]\,e_{a_x}[k] \\
\dxthreehat[k{+}1] &= \lc\,\dxthreehat[k] - d(1{-}\lc)\,\xDhat[k] + \tfrac{d(d{+}1)}{2}(1{-}\lc)\,\dxDhat[k] \\
                   &\quad + (1{-}\lc)\mathcal{F}_2[k]\,\daxhat[k] + \ell_{31}[k]\,e_x[k] + \ell_{32}[k]\,e_{a_x}[k] \\
\xDhat[k{+}1]      &= \xDhat[k] + \dxDhat[k] + \ell_{41}[k]\,e_x[k] + \ell_{42}[k]\,e_{a_x}[k] \\
\dxDhat[k{+}1]     &= \dxDhat[k] + \ell_{51}[k]\,e_x[k] + \ell_{52}[k]\,e_{a_x}[k] \\
\axhat[k{+}1]      &= \axhat[k] + \daxhat[k] + \ell_{61}[k]\,e_x[k] + \ell_{62}[k]\,e_{a_x}[k] \\
\daxhat[k{+}1]     &= \daxhat[k] + \ell_{71}[k]\,e_x[k] + \ell_{72}[k]\,e_{a_x}[k]
\end{aligned}\right.
\label{eq:estimator}
\end{equation}
where $\ell_{ij}[k]$ are the entries of the $7\times 2$ estimator feedback matrix $\LL[k]$. The Row-3 predict terms in $\xDhat$, $\dxDhat$, and $\daxhat$ are chosen so that subtracting \eqref{eq:estimator} from \eqref{eq:state-evo} reproduces the error dynamics \eqref{eq:err-dyn}.

The dynamics of estimation errors $\ee[k] = \xx[k] - \xxhat[k]$ is obtained by subtracting \eqref{eq:estimator} from \eqref{eq:state-evo}:
\begin{equation}
\ee[k{+}1] = \Fe[k]\,\ee[k] - \LL[k]\,\HH\,\ee[k] + \qq[k],
\label{eq:err-dyn}
\end{equation}
where $\qq[k]$ is the process noise and $\Fe[k]$ is the state-transition matrix of the error dynamics given in \S\ref{sec:Fe}.

%=====================================================================
\section{State Transition Matrix}
\label{sec:Fe}
%=====================================================================

By the second-order process for $\xD$ in \eqref{eq:state-evo}, each past value $\xD[k{-}i] = \xD[k] - i\,\dxD[k]$. Summing over $i = 1, \dots, d$ and using $\sum_{i=1}^{d} i = d(d{+}1)/2$,
\begin{equation}
\sum_{i=1}^{d}\xD[k{-}i] = d\,\xD[k] - \tfrac{d(d{+}1)}{2}\,\dxD[k].
\label{eq:xD-sum}
\end{equation}
Substituting \eqref{eq:xD-sum} into Row~3 of \eqref{eq:state-evo} and absorbing the state-dependent part of $\etaR$ (via $\mathcal{F}_1,\mathcal{F}_2$ from \eqref{eq:F12}),
\begin{equation}
\Fe[k] = \begin{bmatrix}
0 & 1 & 0 & 0 & 0 & 0 & 0 \\
0 & 0 & 1 & 0 & 0 & 0 & 0 \\
0 & 0 & \lc & -(1{+}d(1{-}\lc)) & \tfrac{d(d{+}1)}{2}(1{-}\lc) & -\fdx[k]-(1{-}\lc)\mathcal{F}_1[k] & (1{-}\lc)\mathcal{F}_2[k] \\
0 & 0 & 0 & 1 & 1 & 0 & 0 \\
0 & 0 & 0 & 0 & 1 & 0 & 0 \\
0 & 0 & 0 & 0 & 0 & 1 & 1 \\
0 & 0 & 0 & 0 & 0 & 0 & 1
\end{bmatrix}.
\label{eq:Fe}
\end{equation}

%=====================================================================
\appendix
\newpage
\section{Assumptions and Design Alternatives}
\label{app:assumptions}
%=====================================================================

\subsection{Catalog of Forced Approximations}
\label{app:catalog}

\begin{center}
\renewcommand{\arraystretch}{1.2}
\begin{tabular}{@{}llp{0.50\linewidth}@{}}
\toprule
Tag & Substitution & Reason \\
\midrule
F1 & $\ax[k] \mapsto \axhat[k]$ & True $\ax$ is unobservable; only the EKF posterior is available. \\
F2 & $\ax[k{-}i] \mapsto \axhat[k{-}i]$ & Past true $\ax$ is unobservable; buffered past posteriors used. \\
F3 & $x[k{-}d] \mapsto x_m[k]$ (drop $n_x$) & Sensor noise $n_x$ is unobservable. \\
F4 & $\fT[\cdot] \mapsto 0$ & Zero-mean white thermal force; no predictor exists. \\
F5 & $\xD[k] \mapsto \xDhat[k]$ & True $\xD$ is unobservable; only the EKF posterior is available. \\
F6 & $\xD[k{-}i] \mapsto \xDhat[k{-}i]$ & Past true $\xD$ is unobservable; buffered past posteriors used. \\
\bottomrule
\end{tabular}
\end{center}

All six are forced: the substituted quantity is either unobservable (F1, F2, F5, F6) or zero-mean white noise without a predictor (F3, F4). The single design parameter $\lc \in (0,1)$ controls the closed-loop convergence rate versus noise injection trade-off.

\subsection{Legacy Passive Form}
\label{app:passive}

Collapsing the buffered past estimates in \eqref{eq:ctrl-law} to their current values,
\begin{equation}
\axhat[k{-}i] \mapsto \axhat[k], \qquad \xDhat[k{-}i] \mapsto 0 \quad (i \geq 1),
\label{eq:collapse}
\end{equation}
the control law simplifies to
\begin{empheq}[box=\fbox]{align}
\fdx[k] = \,& \axhat^{-1}[k]\bigl\{x_d[k{+}1] - \lc\,x_d[k] - (1{-}\lc)\,x_d[k{-}d] + (1{-}\lc)\,\dxm[k]\bigr\} \notag \\
           & - (1{-}\lc)\!\sum_{i=1}^{d}\!\fdx[k{-}i] - \axhat^{-1}[k]\,\xDhat[k].
\label{eq:ctrl-law-passive}
\end{empheq}
This is the form used in \S\ref{sec:trk}--\S\ref{sec:est}. Under \eqref{eq:ctrl-law-passive} the closed-loop \eqref{eq:closed-loop} carries the raw past disturbance $-(1{-}\lc)\!\sum_i \xD[k{-}i]$, and the $\axhat[k{-}i]$ contributions reduce, via the Taylor expansion $\ax[k{-}i] = \axhat[k] + \ea[k] - i\,\dax[k]$, to the $\etaR$ residuals carrying $\ea\,\mathcal{F}_1$ and $\dax\,\mathcal{F}_2$. The Row-3 predict in \eqref{eq:estimator} therefore retains $-d(1{-}\lc)\xDhat + \tfrac{d(d{+}1)}{2}(1{-}\lc)\dxDhat + (1{-}\lc)\mathcal{F}_2\,\daxhat$ so that subtraction reproduces the same $\Fe[k]$ \eqref{eq:Fe}.

The full form \eqref{eq:ctrl-law}, by contrast, replaces $-(1{-}\lc)\!\sum_i \xD[k{-}i]$ by the estimation-error sum $-(1{-}\lc)\!\sum_i \exD[k{-}i]$ in \eqref{eq:closed-loop}, eliminating raw-disturbance leakage; the same $\Fe[k]$ is recovered because the error-dynamics matrix is invariant under the controller--estimator split of responsibility.

\subsection{Relation to Meng 2023 Eq.~17}
\label{app:meng}

Setting $\xD \equiv 0$, $\xDhat \equiv 0$, and treating $\ax$ as a known calibrated constant $\dt/\gamma$ (i.e.\ skipping F1, F2, F5, F6), the passive form \eqref{eq:ctrl-law-passive} reduces to
\begin{equation}
\fdx[k] = (\gamma/\dt)\bigl\{x_d[k{+}1] - \lc\,x_d[k] - (1{-}\lc)\,x_d[k{-}d] + (1{-}\lc)\,\dxm[k]\bigr\}
         - (1{-}\lc)\!\sum_{i=1}^{d}\!\fdx[k{-}i],
\label{eq:ctrl-law-meng}
\end{equation}
which is Eq.~17 of Meng~2023. The present \eqref{eq:ctrl-law} is a strict generalization that estimates $\ax$ online (F1, F2) and incorporates a position-dependent disturbance compensator (F5, F6).

\end{document}
